\documentclass[12pt]{article}
\usepackage[utf8]{inputenc}
\usepackage[spanish]{babel}
\usepackage{amsmath}
\usepackage{amssymb}
\usepackage{geometry}
\geometry{a4paper, margin=2cm}

\begin{document}

\begin{center}
    \textbf{Escuela 4117 | Curso: 5to | Materia: Termodinámica} \\
    \Large \textbf{EXAMEN DE RECUPERATORIO: TRANSFERENCIA DE CALOR}
\end{center}

\vspace{0.5cm}

\noindent Nombre y Apellido: \hrulefill \quad Fecha: \hrulefill

\vspace{0.8cm}

\section*{Instrucciones}
\begin{itemize}
    \item Resuelva cada ejercicio detallando el procedimiento (fórmulas y despejes).
    \item El uso de calculadora está permitido.
    \item Preste atención a las unidades de medida.
\end{itemize}

\hr

\subsection*{1. Escalas Termométricas y Unidades}

a) Un termómetro marca $10^\circ\text{C}$. ¿Cuál es el valor equivalente en grados Fahrenheit ($^\circ\text{F}$)?

\vspace{3cm}

b) Un componente mecánico alcanza los $68^\circ\text{F}$. ¿A cuántos grados Celsius ($^\circ\text{C}$) equivale esta temperatura?

\vspace{3cm}

c) Si la temperatura de un gas es de $27^\circ\text{C}$, exprese este valor en la escala absoluta Kelvin (K).

\vspace{3cm}

d) Determine cuántas calorías equivalen exactamente a 10 BTU. (Considere que $1\text{ BTU} = 252\text{ cal}$).

\vspace{3cm}

\hr

\subsection*{2. Calorimetría (Energía y Masa)}

Se dispone de una masa de $5\text{ kg}$ de agua. Calcule la cantidad de calor ($Q$) necesaria para incrementar su temperatura en $20^\circ\text{C}$ (por ejemplo, de $20^\circ\text{C}$ a $40^\circ\text{C}$). (Considere el calor específico del agua $C_e = 1\text{ kcal/kg}^\circ\text{C}$).

\vspace{6cm}

\hr

\subsection*{3. Conducción Térmica (Flujo en Placas)}

Calcule el flujo de calor ($Q$) que atraviesa una placa de hormigón con las siguientes características:

\begin{itemize}
    \item Superficie ($A$): $2\text{ m}^2$
    \item Espesor ($L$): $10\text{ cm}$ (Recuerde pasar a metros: $0.1\text{ m}$)
    \item Coeficiente térmico ($k$): $1.5\text{ kcal/h}\cdot\text{m}\cdot^\circ\text{C}$
    \item Diferencia de temperatura ($\Delta T$): $30^\circ\text{C}$
\end{itemize}

\vspace{6cm}

\end{document}